\begin{homeworkProblem}[Seção 4-4]
  (Princípio de Cavalieri) Sejam $X,Y\subseteq \mathbb{R}^{m+1}$ conjuntos J-mensuráveis tais que, para cada $t\in\mathbb{R}$, as seções
  \begin{align*}
    X_{t}=\{x\in\mathbb{R}^{m}:(x,t)\in X\},\quad Y_t=\{y\in\mathbb{R}^{m}:(y,t)\in Y\}
  \end{align*}
  são J-mensuráveis em $\mathbb{R}^{m}$ e satisfazem $Vol(X_t)=Vol(Y_t)$. Então $Vol(X)=Vol(Y)$. 
  \begin{solution}
    Note que como $X_t$ é J-mensurável, temos que se definimos $\xi_{X_t}:\mathbb{R}^{m}\to\mathbb{R}$ dada por
    \begin{align*}
      \xi_{X_t}(x)=\begin{cases}
        1, &\text{ se } (x,t)\in X \text{,} \\
        0, &\text{ no outro caso }.
      \end{cases}
    \end{align*}
    é uma função integrável e cumpre
    \begin{align*}
      \int_{\mathbb{R}^{m}}\xi_{X_{t}}(x)\,dt=Vol(X_t).
    \end{align*}
    Agora, note que se definimos $\xi_{X}:\mathbb{R}^{m}\times\mathbb{R}\to\mathbb{R}$ dada por 
    \begin{align*}
      \xi_{X}(x,t)=\begin{cases}
        1, &\text{ se } (x,t)\in X \text{,} \\
        0, &\text{ no outro caso }.
      \end{cases}
    \end{align*}
    Então dado $t\in\mathbb{R}$ fixo, temos que $\xi_{X}(x,t)=\xi_{X_t}(x)$ e
    \begin{align*}
      \int_{\mathbb{R}}\xi_{X}(x,t)\,dx=Vol(X_t).
    \end{align*}
    Más note que como $X$ é J-mensurável, temos que
    \begin{align*}
      \int_{\mathbb{R}^{m+1}}\xi_{X}(z)\,dz=Vol(X).
    \end{align*}
    Logo pelo anterior e o teorema da integração repetida temos que
    \begin{align*}
      Vol(X)=&\int_{\mathbb{R}^{m+1}}\xi_{X}(z)\,dz\\
      =&\int_{\mathbb{R}}\int_{\mathbb{R}^{m}}\xi_{X}(x,t)\,dx\,dt\\
      =&\int_{\mathbb{R}}\int_{\mathbb{R}^{m}}\xi_{X_t}(x)\,dx\,dt\\
      =&\int_{\mathbb{R}}Vol(X_t)\,dt.
    \end{align*}
    Logo note que podemos fazer o mesmo para $Y_t$, isto é que
    \begin{align*}
      Vol(Y)=&\int_{\mathbb{R}^{m+1}}\xi_{Y}(z)\,dz\\
      =&\int_{\mathbb{R}}\int_{\mathbb{R}^{m}}\xi_{Y}(y,t)\,dy\,dt\\
      =&\int_{\mathbb{R}}\int_{\mathbb{R}^{m}}\xi_{Y_t}(y)\,dy\,dt\\
      =&\int_{\mathbb{R}}Vol(Y_t)\,dt.
    \end{align*}
    Logo temos que como $Vol(X_t)=Vol(Y_t)$ para cada $t\in \mathbb{R}$, então 
    \begin{align*}
      Vol(X)=\int_{\mathbb{R}}Vol(X_{t})\,dt=\int_{\mathbb{R}}Vol(Y_t)\,dt=Vol(Y).
    \end{align*}
    O que conclui a prova.
  \end{solution}
\end{homeworkProblem}
